\section{Related Work}
\label{sec:related}

\subsection{Industrial Process Control Under Model Mismatch}

Industrial process-control systems usually separate tracking, protection, and actuator coordination across several software and hardware layers. For boiler-turbine operation, a new method is most useful when it can be inserted into the existing measurement and control loop without replacing the coordinated controller. This study therefore treats RoCBF-Net as a supervisory safety filter rather than another standalone controller. The evaluation is tied to a specific boiler-turbine loop, with measured variables, process constraints, disturbance scenarios, and computational latency reported explicitly.

This positioning follows the evidence style used in recent applied control studies. Wu et al.~\cite{wu2026pressure} evaluate pressure regulation through delay, disturbance rejection, and response-index improvements in a low-pressure vacuum process. Feng and Li~\cite{feng2026event} study energy-system voltage regulation with attention to event-triggered updates and communication burden. Teke and Tan~\cite{teke2026reference} tune sliding-mode parameters from specified response behavior rather than from purely formal gain choices. These papers motivate the same presentation choices used here: process variables are named, response and violation curves are reported, and the contribution is framed as a deployable control layer with explicit tuning and latency implications.

\subsection{Constrained Control and Safety Filtering}

MPC is the dominant constrained control methodology in process industries~\cite{qin2003survey,mayne2000constrained}, and offset-free variants address steady-state model mismatch~\cite{morari1999model}. Recent boiler-turbine MPC work in the SAGE measurement and control literature still treats nonlinear feasibility, input constraints, and load-response capability as central design challenges~\cite{abdelbaky2025ioflmpc}. Cautious MPC~\cite{hewing2020cautious} incorporates GP uncertainty into the prediction horizon and provides chance-constrained operation at increased computational cost. Here, NMPC is used as a strong benchmark because it directly optimizes over a short horizon. RoCBF-Net instead performs a one-step projection intended for low-latency safety intervention.

Control barrier functions enforce forward invariance through QP-based control projection~\cite{ames2017cbf,wieland2007constructive}. HOCBFs~\cite{xiao2021hocbf} extend this idea to high-relative-degree constraints. Robust CBF formulations~\cite{taylor2020issf,jankovic2018robust} add margins based on structural disturbance bounds. Related nonlinear-control studies continue to emphasize uncertainty, stochastic perturbations, input saturation, and actuator faults as practical barriers to safe operation~\cite{fang2026adaptive,dou2026nonsingular}. These margins are often applied at their full theoretical value. The present results show why this choice can be counterproductive in a process-control QP: $\epsilon_\kappa=1.0$ over-constrains the feasible set in the tested benchmark conditions, whereas smaller calibrated values preserve both safety and control authority.

\subsection{GP-Enhanced CBFs}

GPs provide nonparametric residual learning with uncertainty quantification~\cite{rasmussen2006gp,srinivas2010gp}. In measurement practice, Bayesian machine learning has recently been used to estimate measurement noise and provide uncertainty information from limited calibration data~\cite{alburayt2026bayesian}. RoCBF-Net uses the same uncertainty-aware logic for a control purpose: the GP posterior mean corrects the model residual, while the posterior variance is propagated into a safety margin. Several GP-CBF methods apply posterior uncertainty to first-order CBF constraints~\cite{cheng2019guaranteed,fisac2018general,castaneda2021pointwise}. Aali and Liu~\cite{aali2024learning} extend GP-HOCBF integration to high-order constraints through a second-order cone reformulation, and Zhang et al.~\cite{zhang2025adaptive} develop an adaptive sparse-GP HOCBF filter for robot control. Lederer et al.~\cite{lederer2021uniform} provide uniform GP error bounds that are complementary to this line of work.

RoCBF-Net differs in two ways. First, uncertainty is propagated recursively through the HOCBF $\psi$-chain rather than applied only as a top-level aggregate. Second, the margin is treated as a commissioning parameter. This distinction is central in boiler-turbine operation, where a theoretically valid margin can still make the real-time QP infeasible when actuator authority is limited.

\subsection{Safety Margin Calibration}

The safety-filter literature often treats the robustness margin as a certificate requirement rather than an operating parameter. Industrial commissioning, in contrast, routinely selects gains, thresholds, and protection margins from process tests. This is consistent with recent response-oriented control tuning work, where a desired closed-loop response is used to select controller parameters and then validated through response curves and performance indices~\cite{teke2026reference,wu2026pressure}. This paper brings that practice into GP-HOCBF filtering by introducing $\epsilon_\kappa$ as a tunable factor. The experiments identify three regimes: GP mean correction is sufficient for additive uncertainty; a small positive margin is needed for moderate state-dependent coupling; and beyond the deployment envelope, increasing the margin worsens feasibility and should trigger fallback rather than tighter filtering.

In the broader safe-control literature, shielding approaches~\cite{alshiekh2018safe} correct policy actions, constrained reinforcement learning methods~\cite{achiam2017constrained} enforce expected constraints, and adaptive CBFs~\cite{long2023ecbf} estimate uncertain parameters online. RoCBF-Net is closest to shielding in architecture, but its intervention is model-based and constraint-specific. It returns the nearest action satisfying pressure, enthalpy, and power HOCBF constraints under a calibrated GP residual model.
